% =====================================================================
%  Section 6 -- Discussion
% =====================================================================

\section{Discussion}
\label{sec:discussion}

\subsection{The Negative-KD Finding and What It Means for the Field}
\label{sec:discussion-negative-kd}

The central empirical claim of this paper is that knowledge distillation does not help on FSSD-class flood segmentation when the student is an ImageNet-pretrained lightweight UNet. We frame this as a \emph{negative result} rather than a \emph{null result} because the comparison is not ``we did not find a positive effect'' but rather ``the best student (MobileViT-XXS, no-KD) \emph{exceeds} the teacher (UNet$+$EfficientNet-B0)''. Distillation on this benchmark is at best redundant and at worst counter-productive: combined KD on MobileViT-XXS produced $-0.005$ IoU vs.~the same architecture trained without it.

We interpret this mechanistically. The classical motivation for KD --- Hinton's \emph{dark knowledge} argument~\cite{hinton2015distilling} and subsequent feature-map alignment work~\cite{romero2015fitnets,liu2019structured,wang2020intra,he2021channel} --- assumes that the student has insufficient capacity or pretraining to recover the teacher's behaviour from labelled data alone, so the teacher's soft probabilities and intermediate features provide a richer supervision signal than the binary mask. On FSSD, that assumption is violated in two directions: (i)~the benchmark has only 663 images and saturates at IoU $\approx 0.93$ --- the labelled signal is already nearly sufficient --- and (ii)~the student's ImageNet pretraining already supplies the visual priors (edge detectors, texture filters, semantic-region features) that a teacher trained on the same FSSD would have been transferring. The teacher's ``dark knowledge'' is therefore largely redundant with the student's pretraining, and the response-based KL term reduces to a noise-injection regulariser. The feature-MSE term performs slightly better --- it provides at least \emph{some} signal beyond the labels --- but the gain is well within fold variance and is not statistically significant at our reportable resolution.

This finding has two implications for the flood-segmentation literature:

\begin{enumerate}
\item \textbf{Benchmarks matter.} A KD recipe that wins on Cityscapes (19 classes, 5000 images, IoU $\approx 0.80$--$0.85$) does not transfer to FSSD (2 classes, 663 images, IoU $\approx 0.93$). Future KD work in this space should report headroom (teacher--best-pretrained-student gap) \emph{before} claiming distillation effects.
\item \textbf{Modern ImageNet pretraining is doing the work.} The 2026 wave of \texttt{timm}-pretrained lightweight backbones brings the no-KD baseline so high that compression strategies \emph{other than} distillation --- quantization, pruning, architecture choice --- dominate the deployment-relevant trade-offs. We recommend that authors investing in flood-segmentation compression budget their effort accordingly.
\end{enumerate}

\subsection{Per-Platform Deployment Recommendations}
\label{sec:discussion-deployment}

Based on the joint segmentation $\times$ latency $\times$ stability evidence (Sections~\ref{sec:results-kd}--\ref{sec:results-latency}), we make the following deployment recommendations:

\begin{itemize}
\item \textbf{Field-laptop / tablet inference (Apple Silicon class):} EfficientNet-Lite0 UNet INT8. 17\,ms p50 (\deployableMTwoFPS~FPS) at IoU 0.885, $3.83\times$ size reduction. Stable across folds ($\pm 0.015$). This is the configuration we ship in our public release.
\item \textbf{Browser citizen-science / web-portal inference:} MobileNetV3-Small UNet FP32 (133\,ms p50; 7.5\,FPS), or MobileViT-XXS UNet FP32 (149\,ms p50; 6.7\,FPS) for slightly higher accuracy. INT8 is not yet supported in ONNX Runtime Web 1.17 and is gated on ORT-Web 1.18$+$ shipping a QDQ INT8 Conv kernel.
\item \textbf{Maximum-accuracy off-line inference:} MobileViT-XXS UNet FP32 (28\,ms p50 on M2~Pro CPU; IoU 0.937). Recommended for human-in-the-loop annotation pipelines.
\end{itemize}

We still do not recommend the teacher (UNet$+$EffB0) for deployment, but for a subtler reason than raw FP32 speed. On a common ONNX Runtime CPU backend the teacher runs at 37.6\,ms p50 (\teacherMTwoFPS~FPS), only $1.1$--$1.9\times$ slower than the FP32 students, so the case against it rests on \emph{model size and quantization headroom} rather than FP32 latency: the recommended EfficientNet-Lite0 INT8 is $2.2\times$ faster (17.3 vs 37.6\,ms) and roughly $4\times$ smaller (5.2 vs 22.3\,MB) while recovering 95\% of the teacher's IoU, and the teacher does not INT8-quantize within our budget.

Table~\ref{tab:t7-decision} consolidates these recommendations into a decision matrix, pairing each operational scenario with the configuration we recommend and the metric that decides it.

% =====================================================================
%  Table T7 -- Deployment Decision Matrix
%  Source of truth: summary.json (FP32 IoU), quantized.json (INT8 IoU),
%  lat_m2.json / lat_wasm.json (latency), T1 (size), T6 (OOD).
%  Re-presentation of the Section 6.2 recommendations; no new experiments.
% =====================================================================
\begin{table*}[ht]
\centering
\caption{\textbf{Deployment decision matrix.} Recommended configuration per operational scenario, with the deciding metrics. All values are drawn from Tables~\ref{tab:t1-footprints}, \ref{tab:t2-segmentation}, \ref{tab:t4-quant}, \ref{tab:t5-latency}, and~\ref{tab:t6-ood}; this table re-presents the Section~\ref{sec:discussion-deployment} recommendations as an at-a-glance practitioner reference and introduces no new experiments. Latency is p50 on the platform relevant to each scenario (Apple~M2~Pro ORT-CPU, or browser WASM SIMD).}
\label{tab:t7-decision}
\small
\setlength{\tabcolsep}{4pt}
\begin{tabular}{p{2.9cm}lcccp{4.5cm}}
\toprule
\textbf{Deployment scenario} & \textbf{Recommended} & \textbf{IoU} & \textbf{Latency (p50)} & \textbf{Size (MB)} & \textbf{Rationale / caveat} \\
\midrule
Field laptop / tablet (Apple-Silicon CPU)
  & \textbf{EffLite0-UNet INT8} & 0.885 & 17\,ms (M2) & 5.2
  & Fastest deployable operating point; the only student that INT8-quantizes stably ($\pm 0.015$). Shipped in our release. \\
\addlinespace
Browser / citizen-science web portal
  & MNV3-S-UNet FP32 & 0.927 & 133\,ms (WASM) & 13.8
  & Fastest under single-threaded WASM SIMD; MobileViT-XXS FP32 (149\,ms, IoU 0.937) if accuracy matters more. No INT8 in ORT-Web~1.17. \\
\addlinespace
Maximum-accuracy off-line / human-in-the-loop
  & MobileViT-XXS-UNet FP32 & 0.937 & 28\,ms (M2) & 12.0
  & Highest IoU of any configuration; FP32 only (its INT8 collapses, Table~\ref{tab:t4-quant}). \\
\addlinespace
Cross-sensor / satellite optical
  & \emph{retrain / domain-adapt} & 0.16\textsuperscript{$\dagger$} & --- & ---
  & No FSSD-trained model transfers zero-shot (Table~\ref{tab:t6-ood}); requires domain adaptation or in-sensor training data. \\
\bottomrule
\end{tabular}
\\[3pt]
{\footnotesize \textsuperscript{$\dagger$}Zero-shot Sen1Floods11 IoU; all models $\approx 0.16$ (Section~\ref{sec:results-ood}).}
\end{table*}


\subsection{Architecture-Dependent PTQ Stability is a Practitioner Trap}
\label{sec:discussion-ptq-trap}

The 0.45-IoU PTQ collapse on MobileNetV3-Small and MobileViT-XXS would not be predicted from the classification-quantization literature, which typically reports $\sim 1$ IoU point PTQ tax on ResNet/MobileNet/EfficientNet families on ImageNet~\cite{krishnamoorthi2018quantizing,jacob2018quantization}. The flood-segmentation regime differs in two respects: (i)~the output is a per-pixel binary mask, so any quantization error that biases the post-sigmoid threshold materially shifts the segmentation; and (ii)~flood images contain a wide dynamic range of bright reflective surfaces that drive calibration statistics in HardSwish-based backbones. We have observed that the same Percentile-99.999\% recipe that fully recovers EfficientNet-Lite0 drives MobileNetV3-Small to sharply different per-fold INT8 outcomes: across folds 0--2 its INT8 IoU is $0.772$, $0.0001$ (near-total collapse to the background class), and $0.657$~--- one catastrophic failure and two partial degradations, with no fold recovering to its FP32 level ($\approx 0.93$). The large cross-fold standard deviation ($\pm 0.34$) reflects this instability rather than a smooth degradation, which is why the per-fold values, not the mean, are the informative summary here.

The practical guidance: when assembling a flood-segmentation deployment artifact under a static-QDQ-INT8 budget, prefer backbones with bounded activation ranges (ReLU6 is the safest, ReLU is fine, HardSwish is risky, GELU and softmax-attention are \emph{not} PTQ-friendly without QAT). Without this guidance, a practitioner running the standard ONNX Runtime quantization tutorial on MobileNetV3-Small would conclude --- incorrectly --- that the model ``doesn't quantize'' rather than that the recipe is mismatched to the activation family.

\subsection{Failure Modes (Qualitative)}
\label{sec:discussion-failures}

Figure~\ref{fig:f3-qualitative} shows that both teacher and EfficientNet-Lite0-INT8 student share two persistent failure modes:
\begin{itemize}
\item \textbf{Bright-reflective-surface confusion:} wet asphalt, glare-on-roof, and large reflective puddles are misclassified as flood. This is a labelling-modality issue, not a model-capacity issue --- RGB alone cannot disambiguate water from wet non-water surfaces.
\item \textbf{Vegetation-occluded edges:} when flooded water is partially overhung by trees or tall grass, both models under-segment the flood boundary. The 5-pixel margin at the boundary contributes disproportionately to the residual IoU gap.
\end{itemize}

Both failure modes argue for SAR or multispectral fusion rather than for a different RGB segmentation recipe. We discuss this in Section~\ref{sec:limitations}.

Table~\ref{tab:t8-taxonomy} organises the failure modes observed across the whole study --- the two RGB-modality failures above, the cross-sensor distribution shift of Section~\ref{sec:results-ood}, and the INT8 PTQ instability of Section~\ref{sec:results-quant} --- into a single taxonomy, noting for each its mechanism, error direction, the models affected, and the mitigation it implies. The categories are qualitative: we report the error direction (which matters for the operational cost asymmetry of Section~\ref{sec:discussion-asymmetry}) but do not claim per-class frequencies.

% =====================================================================
%  Table T8 -- Error Taxonomy
%  Qualitative synthesis of failure modes established in
%  Sections 5.3 (PTQ), 5.7 (OOD), 6.4 (qualitative failure modes).
%  No new experiments; per-class frequencies are NOT claimed.
% =====================================================================
\begin{table*}[ht]
\centering
\caption{\textbf{Error taxonomy for FSSD flood segmentation.} Four failure classes observed across this study, their mechanism, the direction of the resulting error (operationally, false negatives are the costlier direction, Section~\ref{sec:discussion-asymmetry}), the models affected, and the mitigation each implies. Categories are qualitative, synthesised from Sections~\ref{sec:results-quant}, \ref{sec:results-ood}, and~\ref{sec:discussion-failures}; we do not claim per-class frequencies.}
\label{tab:t8-taxonomy}
\small
\setlength{\tabcolsep}{4pt}
\begin{tabular}{p{2.6cm}p{4.2cm}p{2.1cm}p{2.4cm}p{3.2cm}}
\toprule
\textbf{Failure class} & \textbf{Mechanism} & \textbf{Error type} & \textbf{Affected models} & \textbf{Mitigation} \\
\midrule
Reflective-surface confusion
  & RGB cannot disambiguate water from wet non-water surfaces (wet asphalt, roof glare, standing puddles)
  & False positive
  & Teacher and all students
  & SAR / multispectral fusion (outside the RGB scope) \\
\addlinespace
Vegetation-occluded boundary
  & Flood edges overhung by trees or tall grass are under-segmented
  & False negative
  & Teacher and all students
  & Modality limitation; boundary-aware loss at best partial \\
\addlinespace
Cross-sensor distribution shift
  & Close-range RGB $\rightarrow$ Sentinel-2 satellite optical domain gap
  & Whole-image collapse (IoU\,$\approx$\,0.16)
  & All (uniformly)
  & Domain adaptation or in-sensor retraining \\
\addlinespace
INT8 PTQ instability
  & Unbounded HardSwish tail / attention dynamic range exceed static-QDQ calibration
  & Model-level collapse ($\Delta$IoU\,$\approx-0.45$, high cross-fold variance)
  & MobileNetV3-Small, MobileViT-XXS
  & QAT or a ReLU6-bounded backbone (EffLite0 unaffected) \\
\bottomrule
\end{tabular}
\end{table*}


\subsection{Operational Error Asymmetry}
\label{sec:discussion-asymmetry}

The metrics we report --- IoU and F1 --- weight false positives and false negatives equally, but in flood response the two errors carry very different operational costs. A false negative (unmapped flood) can translate into under-warning and under-allocation of rescue resources, whereas a false positive (over-flagged flood) mainly costs excess caution; for most disaster-response decisions the false negative is the more expensive error. We report precision and recall separately in Table~\ref{tab:t2-segmentation} precisely so this trade-off is visible rather than hidden inside a single symmetric score: the FP32 teacher is close to balanced (precision $0.962$, recall $0.960$), and the two qualitative failure modes above pull in opposite directions --- bright-surface confusion inflates false positives, vegetation occlusion inflates false negatives.

INT8 conversion further perturbs this balance, and not in a stable direction: the deployable EfficientNet-Lite0 student's per-fold (precision, recall) after quantization are $(0.965, 0.934)$, $(0.917, 0.963)$, and $(0.897, 0.963)$ --- recall-favouring (safer, over-segmenting) in two folds but precision-favouring (missing flood) in the third. We therefore make two recommendations for operational use. First, tune the post-sigmoid decision threshold toward recall --- accepting more false positives to reduce missed flood --- rather than leaving it at the default $0.5$, matching the asymmetric cost structure. Second, re-calibrate that threshold \emph{after} INT8 conversion, since quantization shifts the operating point by an amount and in a direction our results show is not predictable from the FP32 model alone. A full cost-weighted evaluation (e.g.\ an $F_\beta$ score with $\beta>1$, or a region-level warning-recall metric) is left to future work: the per-image cost model is application- and jurisdiction-specific and cannot be fixed once for all deployments.
